\documentclass[../main.tex]{subfiles}
\begin{document}

\section{Conclusions}
\label{sec:conclusions}

Eight state-of-the-art LLMs assessed RoB~1.0 for 227 randomised controlled trials
in a zero-shot, single-agent setting.
Overall accuracy ranged narrowly from 57.6\% to 62.5\%, but Macro-F1 and Gwet's~AC1
revealed meaningful differences in how models handled the three nominal classes.
Models with the highest accuracy achieved it partly by concentrating predictions on
majority-class labels, a strategy that inflates accuracy whilst suppressing detection
of High Risk and Unclear Risk judgements---the very categories most consequential for
evidence appraisal.
Macro-F1 and Gwet's~AC1 are therefore more informative primary metrics than accuracy
for this task.

\emph{Unclear Risk} was the most consistently under-detected class across all models
and domains, with recall below 0.20 in three of seven domains.
This systematic failure reflects the epistemic character of the category---it
requires the model to recognise that key methodological information is absent,
a more complex inference than identifying a positive or negative methodological
signal in the text.

Domain-level and stratification analyses revealed structured heterogeneity in
performance.
Sequence generation~(D1) and allocation concealment~(D2) were reliably assessed
across all models, whereas other sources of bias~(D7) showed the greatest
inter-model accuracy spread and blinding of outcome assessors~(D4) had the lowest
oracle ceiling (78.0\%).
Clinical domains with clear, domain-specific methodological criteria (Orthopaedics
\& Trauma, Endocrine \& Metabolic) supported higher Macro-F1 than domains with more
heterogeneous study designs (Infectious Disease, Reproductive \& Sexual Health).

The oracle accuracy of 85.8\%---23 percentage points above the best individual
model---confirms that the eight models make substantially non-redundant errors,
establishing a firm empirical ceiling for any multi-model selection strategy.

The comparison between per-domain and single-call inference modes
(Section~\ref{sec:inference}) produced small differences that do not
collectively support a performance advantage for either strategy.
Ten of twelve $\Delta$ values across both models and six metrics had 95\%~CIs
spanning zero; the two exceptions were class-specific effects of limited
practical magnitude (Unclear Risk F1 for Flash: $-$0.044; Low Risk F1 for
DeepSeek: $+$0.026) that did not propagate to aggregate metrics.
Most notably, D7~(Other Bias)---the domain for which joint contextual
assessment was hypothesised to be most beneficial---showed negligible
differences under both modes ($\pm$1~pp), indicating that its difficulty
is intrinsic to the vagueness of the criterion rather than a consequence
of fragmented inference calls.
Taken together, the preliminary evidence from this two-model comparison does
not identify a meaningful performance rationale for preferring single-call
over per-domain assessment; the primary practical distinction between the
two modes is computational (single-call reduces inference calls by 7$\times$)
rather than one of accuracy.

Several limitations should be noted.
First, the ground-truth annotations derive from published Cochrane review
reports and were not independently re-annotated for this study.
The Cochrane annotations represent the consensus output of a structured
multi-reviewer process (independent assessment by at least two trained
reviewers, with disagreements resolved by discussion or arbitration), which
is the reference standard used operationally by the field and the most
appropriate benchmark for evaluating a tool intended to assist real-world
systematic review workflows.
Quantifying internal reliability within this corpus would require re-annotating
a subsample with independent trained reviewers; however, agreement between
Cochrane-trained reviewers and non-specialist re-annotators would conflate
task difficulty with expertise asymmetry---the same confound that limits the
interpretability of comparisons between Cochrane and external reviewers
\citep{ArmijoOlivo2014}.
The relevant population-level IRR estimate is therefore taken from Hartling
et al.\ \citep{Hartling2013}, who measured agreement between formally trained
reviewer pairs applying the same protocol: $\kappa = 0.79$ for D1 but only
$\kappa = 0.24$--$0.37$ for remaining domains, falling to $\kappa = 0.05$--$0.09$
between independent consensus pairs; independent Cochrane teams show
similarly variable agreement \citep{Konsgen2020}.
These figures establish that the reference standard itself carries substantial
noise, placing a practical ceiling on any automated system's agreement with
published annotations that is independent of model quality.
Converting the domain-level $\kappa$ values for the non-D1 domains
(0.24--0.37) to expected accuracy using the observed pooled class distribution
(LR\,53.5\%, HR\,17.4\%, UR\,29.1\%) yields approximately 58--62\%---a range
that coincides closely with the model accuracy range observed in this study
(57.6\%--62.5\%), suggesting that the evaluated LLMs perform at a level broadly
comparable to a second independent human reviewer across most domains.
Second, each model was queried once per study using provider-default generation
parameters.
This reflects the intended evaluation scope: zero-shot performance as models
are available to practitioners, without artificial parameter tuning.
Manipulating generation temperature would introduce heterogeneous conditions
across model families, and strategies such as majority voting over repeated
draws represent post-hoc optimisation outside the zero-shot benchmark design.
As a consequence, differences between models separated by fewer than five
percentage points should not be interpreted as definitive rank orderings in the
absence of confidence intervals; these are discussed as directional trends.
Third, all models were accessed in June~2026 via the OpenRouter~API using the
identifier strings reported in Table~\ref{tab:models}; providers may silently
update the weights associated with a fixed identifier, so exact reproduction
requires accessing the same model strings within the same inference window.
Results should therefore be interpreted as reflecting the checkpoint state active
in June~2026, not any future revision of the same identifier.
Fourth, the single-agent zero-shot design examines one specific slice of the
capability space; few-shot prompting, retrieval augmentation, fine-tuning, and
multi-agent architectures were not evaluated and may yield substantially
different results.
Fifth, training-data contamination cannot be excluded, but several structural
features of the evaluation design limit its plausibility as a primary driver
of results.
Memorisation-based performance would require a model to (i)~have seen the
particular one of the 127 distinct Cochrane~2024 reviews from which each study
was drawn, (ii)~correctly identify that the input Markdown document corresponds
to a specific entry in that review's RoB table, and
(iii)~reproduce the reviewer's label rather than reason from the text.
The Cochrane reviews sampled here were published in~2024; most included RCTs
appear in paywalled journals, making full-text ingestion at training time less
likely for the majority of papers.
Moreover, the empirical pattern is inconsistent with memorisation at scale:
the best individual model (Claude~Haiku~4.5) achieved a perfect 7/7 score on
only 12 of 227 studies (5.3\%), and even the oracle across all eight models
reached 7/7 on only 40 studies (17.6\%).
Both figures are consistent with the rates expected under independent domain
judgements at the observed per-domain accuracy ($0.60^7 \approx 2.8\%$; inflated
modestly by the low intraclass correlation of $\hat\rho = 0.044$ to roughly
4--5\% per model; see Supporting Information~S6 for derivation), leaving no systematic excess attributable to retrieval.
Targeted contamination probes (e.g.\ membership inference or verbatim
reproduction tests) were not conducted, as the opaque training corpora of the
evaluated models preclude rigorous analysis; this remains a caveat for
interpretation.
Sixth, the Cochrane corpus used here comprises better-reported trials than the
broader trial literature, so these performance estimates may be optimistic
relative to unselected primary research.\end{document}
